\section{An actual cubic deformation ring at inert coefficient primes}
\label{sec:cubic-universal-model}

The original named \texttt{Dat} record fixes the rational
character $\Psi$.  The underlying framed representation theorem
has no such restriction.  This section constructs the modified
ring itself, and transports the entire local component at inert
coefficient primes.  It also supplies a homogeneous presentation;
the additional numerical relation count in the source's
\texttt{Model} remains separate.

Keep the original datum and all its \texttt{GD} conditions, let
$K=\mathbb Q(\omega)$, and fix $j\in\{1,2\}$.
Write $\Gamma=G_{\mathbb Q}$, $H=G_K=\ker\eta$,
$\gamma=\kappa_\ell^e\psi$, and choose the normalized
complex conjugation $c$ with $\gamma(c)=1$.
The coefficient DVR $O$ contains $\mu_3$.  If needed at
$\ell\ne3$, adjoining $\mu_3$ is an unramified coefficient
extension, so the source's restriction $e_O=1$, when imposed,
is preserved.  The old local datum is then formed over this
coefficient ring, not obtained by silently assuming an
unchanged residue field.
Define the crossed character and the cubic basepoint by
\[
 \begin{split}
 q(hc^a)&=\chi(h),\qquad
 q(gg')=q(g)q(g')^{\eta(g)},\\
 r_\chi(g).\mathrm{mat}
   &=\operatorname{diag}\bigl(q(g)^j\gamma(g),1,
                               q(g)^j\gamma(g)^{-1}\bigr),\qquad
 r_\chi(g).\mathrm{mult}=\eta(g).
 \end{split}                                                    \tag{G49.1}
\]
The anticyclotomic property of $\chi=\chi_{269,3}$ makes
$r_\chi$ a homomorphism, as checked in AC-K29-kill.
Write $\overline r_\chi$ for its reduction and set
\[
 \mathcal K_\chi=
 \{g\in I_v:\ v\ne\ell,\ \overline r_\chi(g)=1\},\qquad
 M_\chi=\overline{\langle\mathcal K_\chi\rangle}^{\rm normal}.
                                                                  \tag{G49.2}
\]
The union over finite rational places is intended in this notation.

\begin{proposition}[AC-G49-inert: transport of the whole local component]
\label{AC-G49-inert}
Suppose $\ell\ne3,269$ is inert in $K$.
Then $\chi|_{G_{K_\ell}}=1$.  There is a diagonal matrix
\[
             C=\operatorname{diag}(a,1,a),\qquad a\in\mu_3,
                                                                  \tag{G49.3}
\]
for which the restriction of $r_\chi$ to $G_{\mathbb Q_\ell}$
is $\operatorname{conj}_C r_0(\eta,\gamma)$.
Conjugation by $C$ is an equivalence of the entire old and
cubic local framed deformation functors over every complete
Noetherian local $O$-algebra with residue field $k_O$.

Consequently, if $U_0$ is the original local universal ring,
one may use its same underlying ring with universal
representation $\operatorname{conj}_C U_0.r$ as the cubic
local universal object $U_\chi$.  Put
\[
                   P_\chi=I_t(U_0)\subset U_\chi.R.             \tag{G49.4}
\]
This is an actual selected local prime through the cubic
basepoint.  All fields of the source's \texttt{LocGood}
statement, with basepoint replaced by $r_\chi$, hold:
strict conjugation stability and the twelve-dimensional
weighted formal chart, with unchanged multiplicities
\[
 \bigl(d_{-2},d_{-1},d_0,d_1,d_2\bigr)=(1,2,3,4,2).
                                                                  \tag{G49.5}
\]
These conclusions retain the original local published
hypotheses.  The same conclusion holds when $\ell$ splits
in $K$ and $\chi$ is trivial on the local decomposition group.
\end{proposition}
\begin{proof}
\proofstep{The anticyclotomic local character.}
At an inert prime different from $3,269$, $K_\ell/\mathbb Q_\ell$
is unramified quadratic and $\chi$ is unramified.
An element of the rational decomposition group outside $G_K$
conjugates $\chi$ to $\chi^{-1}$.  On the unramified quotient
of $G_{K_\ell}$ this conjugation is trivial: that quotient is
procyclic, generated by the square of rational Frobenius.
Hence $\chi=\chi^{-1}$ on this local group.  Since $\chi^3=1$,
it follows that $\chi=1$ there.

It follows that the restriction of $q$ to the rational local
group is $1$ on $\ker\eta$ and has one value $b\in\mu_3$
on its other coset.  Choose $a=b^{2j}$, so that $a^2=b^j$.
The source involution $\theta(A)=JA^{-t}J$ satisfies
$\theta(C)=C^{-1}$.  Since $C$ commutes with every diagonal
matrix, semidirect conjugation gives
\[
 \bigl(\operatorname{conj}_C r_0(g)\bigr).\mathrm{mat}
   = C\,r_0(g).\mathrm{mat}\,
                 \theta_{\eta(g)}(C^{-1})
   =\begin{cases}
       r_0(g).\mathrm{mat},&\eta(g)=1,\\
       C^2r_0(g).\mathrm{mat},&\eta(g)=-1.
     \end{cases}                                                \tag{G49.6}
\]
This is precisely (G49.1).  For a split prime with trivial
local $\chi$, $q=1$ on the local group and $C=1$ suffices.

\proofstep{The whole functor and its selected prime.}
The two inverse operations $r\mapsto\operatorname{conj}_C r$
and $r'\mapsto\operatorname{conj}_{C^{-1}}r'$ preserve
continuity, invertibility and the fixed multiplier; they
identify the two specified residual matrices.  They commute
with every coefficient-ring homomorphism.  Thus they are
inverse natural transformations of the framed lift functors,
not merely maps between Type II points.  Universality proves
the asserted universal object with the same underlying ring.
Its basepoint map is the original $x_0:U_0.R\to O$, since
its specialized representation is now (G49.6).  Therefore
(G49.4) is prime and contains no nonzero element of $O$,
and the original basepoint lies on it.

For strict conjugation stability, conjugate a strict matrix
$h\equiv1$ by $C$.  Its conjugate is still strict, because
$C$ and $C^{-1}$ are integral.  Apply the original
\texttt{rtConj} to $C^{-1}hC$, and then conjugate back.
This proves the cubic version for every source test ring.

For the chart, use exactly the original map
$\phi:U_0.R\to\operatorname{Frac}(O)[[T_1,\ldots,T_{12}]]$.
Every matrix entry of the conjugated universal representation
is the corresponding original entry times a fixed unit:
the factor is $C_{aa}C_{bb}^{-1}$ when $\eta(g)=1$ and
$C_{aa}C_{bb}$ when $\eta(g)=-1$, since
$\theta(C^{-1})=C$.
Multiplication by these weight-zero units preserves the
weight of every entry.  The kernel, constant coefficient,
linear-coordinate spanning and cotangent injectivity fields
of \texttt{LocChart} are unchanged because the ring map
$\phi$ and its basepoint map are unchanged.  The variable
weights therefore have precisely (G49.5).

The exact source locators are
\path{PadicHodgeTheory/LocGood.lean}, \texttt{locGood},
\path{PadicHodgeTheory/RtChart.lean}, \texttt{rtChart}, and
\path{UnitaryThreeGroups/ConjStable.lean},
\texttt{LocChart}, \texttt{LocGood} and \texttt{locDims}.
Only the explicitly checked replacement of the residual
matrix and basepoint is made; no new literal \texttt{Dat}
instance is asserted.
\end{proof}

\begin{theorem}[AC-G49: existence of the actual normalized cubic ring]
\label{AC-G49}
Retain \texttt{CohInp}, all old \texttt{GD} conditions,
$\ell\ne3,269$, and the local condition (G49.4), supplied by
AC-G49-inert in its stated cases.  Fix the old marked element
$g_1$.  There exists a complete Noetherian local $O$-algebra
$R_\chi$ with residue field $k_O$, an actual $O$-valued
basepoint, and a universal continuous homomorphism
\[
                       r^{\rm univ}:\Gamma\longrightarrow GG_3(R_\chi)
                                                                  \tag{G49.7}
\]
for the following framed functor.  For every test ring $A$
with the source's \texttt{IsCO O A}, its points are exactly
continuous lifts of $\overline r_\chi$ with multiplier $\eta$
which kill $\mathcal K_\chi$, satisfy
\texttt{Normalised}$(c,g_1)$, and whose actual restricted
local universal map kills $P_\chi$.

More generally, the same global construction works for any
actual local unrestricted cubic universal ring and any
ideal whose basepoint map is zero on that ideal.
The construction does not need a global equivalence with
the original rational-seed deformation functor.
\end{theorem}
\begin{proof}
\proofstep{The basepoint satisfies the actual killing condition.}
Take $g\in\mathcal K_\chi$ away from $\ell$.
Its residual multiplier gives $\eta(g)=1$.
The first and third diagonal residual entries give
$\overline q(g)^{2j}=1$.  Together with $q^3=1$ and
$\ell\ne3$, this implies $q(g)=1$.
They then give $\overline\gamma(g)=1$.
The cyclotomic character is unramified at this place, so
$\gamma(g)=\psi(g)$.  The source condition that $\ell$
does not divide the order of $\psi$ makes reduction
injective on its values.  Thus $\gamma(g)=1$ and
$r_\chi(g)=1$.
Continuity shows that $r_\chi$ kills $M_\chi$ as well.
It is normalized because $r_\chi(c).\mathrm{mat}=I$
and its two specified mixed \texttt{zmat} entries vanish.
Its local basepoint map kills (G49.4).

\proofstep{Finiteness for the actual quotient Galois group.}
Let $S$ contain the rational primes dividing
$\ell N D_K\,269$.  Away from $S$, the entire inertia
group lies in $\mathcal K_\chi$.  Therefore
$\Gamma/M_\chi$ is a quotient of $G_{\mathbb Q,S}$.
The source \texttt{GlobalFiniteness} field of
\texttt{CohInp} supplies its \texttt{PhiCond}: for every
open subgroup the continuous homomorphisms to
$\mathbb F_\ell$ are finite.  One way to see the implication
is to apply finite $H^1$ to the finite permutation module
on the cosets of the open subgroup; its cocycles contain
these homomorphisms by the explicit induced-cocycle map.
Pullback then makes the same set finite for every quotient.
This is exactly the proof of
\texttt{phi\_condition\_of\_finite\_cohomology},
\texttt{phi\_surj} and \texttt{phi\_unr} in
\path{UnitaryThreeGroups/RingEx.lean}.
The additional prime $269$ changes only the finite set $S$.
Locally, the unrestricted framed theorem uses
\path{PadicHodgeTheory/PhiCond.lean}, \texttt{phiCond}.

\proofstep{Apply the genuinely general framed theorem.}
The exact theorem
\path{CommutativeAlgebra/UnivLift.lean},
\texttt{UpperTheory.V.Alg.univ\_lift}, takes an
\emph{arbitrary} continuous residual matrix representation.
Apply it to the six-dimensional representation
$\operatorname{emb}\overline r_\chi$ of $\Gamma/M_\chi$.
It gives a complete Noetherian framed ring $R_6$ with the
original residue field.  Impose the following polynomial
relations on its universal matrices $X(g)$.
For $\eta(g)=1$, require
\[
 X(g)=\begin{pmatrix}P&0\\0&Q\end{pmatrix},
                  \qquad P JQ^tJ=I;
                                                                  \tag{G49.8}
\]
for $\eta(g)=-1$, require
\[
 X(g)=\begin{pmatrix}0&P\\Q&0\end{pmatrix},
                  \qquad P JQ^tJ=I.                            \tag{G49.9}
\]
Since the coefficient rings are commutative, the last
identity makes $P$ invertible and forces $Q=JP^{-t}J$.
Thus each matrix is uniquely the embedding of a $GG_3$
element with the prescribed multiplier.  Faithfulness of
this embedding, together with the matrix representation
law, gives the $GG_3$ representation law.  Matrix entries
recover all its entries, so continuity and reduction are
also exact.  Conversely every such $GG_3$ lift satisfies
these equations.  Quotienting by their generated ideal
therefore represents precisely the cubic $GG_3$ lifts
killing $M_\chi$.  The ideal is proper: the actual cubic
basepoint satisfies it.  This rederives, for the new seed,
the complete polynomial \texttt{eqs\_iff}/\texttt{transfer}
construction in \path{UnitaryThreeGroups/UnivEx.lean}.
There is no assumption of absolute residual irreducibility
because the functor remains framed.

\proofstep{Impose the specified local component and normalization.}
Call the last ring $R_\Gamma$.  Restricting its universal
representation gives by local framed universality a unique
map $\iota_R:U_\chi.R\to R_\Gamma$.
Form the ideal generated by $\iota_R(P_\chi)$, by all entries
of $r^{\rm univ}(c).\mathrm{mat}-I$, and by the two entries
\[
                   z(c,g_1)_{01},\qquad z(c,g_1)_{10}.          \tag{G49.10}
\]
Its quotient is $R_\chi$.
The actual basepoint from the first step kills every generator,
so this ideal is proper and the quotient still has
\texttt{IsCO O} and an $O$-valued point.

Every lift satisfying the stated conditions first gives a
unique map from $R_\Gamma$.  The uniqueness of its local
restriction map shows that it kills $\iota_R(P_\chi)$
exactly when its local map kills $P_\chi$.
The remaining equations say exactly \texttt{Normalised}.
The map consequently factors uniquely through $R_\chi$.
The reverse implication holds by evaluating the equations.
This proves the claimed universal property on every source
test ring, including nonreduced rings.
The proof is the explicit quotient argument
\texttt{univ\_quot}, \texttt{univ\_box} and
\texttt{univ\_norm} in
\path{UnitaryThreeGroups/UnivQuot.lean}; their argument uses
only universality and a basepoint satisfying the equations.
\end{proof}

\begin{corollary}[AC-G49-grading: a genuine finite weighted presentation]
\label{AC-G49-grading}
In the setting of AC-G49-inert, the ring $R_\chi$ admits a
finite weighted power-series presentation
\[
  O[[X_1,\ldots,X_s]]\twoheadrightarrow R_\chi,
       \qquad w_i\in\{-2,-1,0,1,2\},                         \tag{G49.11}
\]
whose kernel is homogeneous in the precise sense of
\texttt{WPres.homog}.  Each $X_i$ can be chosen to map to
a matrix-entry difference from the cubic basepoint in the
six-dimensional embedding.  Every such entry difference
is homogeneous of its corresponding weight, and the
basepoint is the origin of this presentation.

This is a weighted presentation of the actual ring.
It does not yet supply the equation counts, smooth ambient
model, derivative lifting, or Selmer comparison fields of
the source's full \texttt{Model} record.
\end{corollary}
\begin{proof}
Give the two copies of the standard three-dimensional
space the same weights $(1,0,-1)$, denoted by $w_6$.
For $s=(g,a,b)$ put
\[
 e_s=\operatorname{emb}(r^{\rm univ}(g))_{ab}
       -\operatorname{emb}(r_\chi(g))_{ab},\qquad
                    w(s)=w_6(a)-w_6(b).                     \tag{G49.12}
\]
These elements lie in the maximal ideal.  They jointly
determine a ring homomorphism, because the embedded
representation determines the original framed lift and
its representing map is unique.

For any target test ring and any unit $t\equiv1$ modulo
its maximal ideal, conjugate the specialized representation
by $T=\operatorname{diag}(t,1,t^{-1})$.
It has the same residual matrix, kills the same elements,
and remains normalized: $\theta(T)=T$, so the matrix at
$c$ stays $I$, and each of the two zero entries in
(G49.10) is only multiplied by a unit.
Its local condition is preserved.  Indeed $C$ commutes
with $T$, and after conjugating by $C^{-1}$ this is the
original strict torus conjugation preserving $I_t$.
Universality therefore supplies another map from $R_\chi$.

The embedded matrix of $T$ is the diagonal matrix with
entries $t^{w_6(a)}$.  Thus conjugation multiplies the
$(a,b)$ entry by $t^{w(s)}$.  The embedded cubic basepoint
has no nonzero entry joining unequal weights: its blocks
are diagonal, including the cross-block form when
$\eta(g)=-1$.  Consequently the transformation is
\emph{exactly}
\[
                           e_s\longmapsto t^{w(s)}e_s.        \tag{G49.13}
\]
This verifies every hypothesis of the literal, seed-independent
\path{CommutativeAlgebra/CoordGrading.lean},
\texttt{UpperTheory.V.Alg.coord\_grading}:
maximal-ideal coordinates, uniqueness from those coordinates,
and their prescribed torus transformation over every test
ring.  That theorem selects finitely many of the $e_s$,
constructs the surjective power-series map, and proves its
kernel homogeneous.  Its weights are the differences of
$(1,0,-1)$, hence are in the stated range.
The basepoint kills all these coordinates, so the composite
with its $O$-valued map is constant-term evaluation.
This proves every assertion of (G49.11).
\end{proof}

\paragraph{Exact remaining model gate.}
The source \texttt{BoxPres} in
\path{UnitaryThreeGroups/IsLocPres.lean} additionally requires,
weight by weight, that the number of relations equals the
number of extra variables plus \texttt{boxDef}.
The theorem \texttt{model\_of\_pres} in
\path{UnitaryThreeGroups/ModelOfPres.lean} uses that equality
and the local chart to construct its specific smooth
ambient model.  No such equality follows just from
(G49.11).  Its cubic verification, the actual mixed tangent
image, and Type I automorphy remain separate steps.
At a split coefficient prime with nontrivial unramified
local $\chi$, even the local chart is not obtained by
(G49.6): the endpoint Frobenius determinant is $\chi^2$.
An assertion of scalar equivalence with the old local
framed functor there would be unjustified.

\paragraph{Geometric dictionary.}
The inert-prime calculation says that the cubic rank-one
local system on $\operatorname{Spec}K$ becomes trivial on
the local decomposition group, while its rational crossed
character is a change of frame.  For finite reductions,
this uses the already audited finite-cover dictionary of
SGA~1, Expos\'e~V, \S\S4 and~7.  The complete local rings
above represent framed lifts with explicit closed conditions;
their existence and selected crystalline component come
from the cited deformation theorems and quotient proofs,
not from that categorical dictionary alone.
